% theme: everyday_functions
\MajorQuestionLesson{身の回りのいろいろな関数}
\SetRowSpace{4mm}

% 制動距離・停止距離
\TypeQuestion{}{everyday_braking_distance}
\TextTall{自動車の速さを $x\mathrm{km/h}$、ブレーキがきき始めてから停止するまでに進む距離（制動距離）を $y\mathrm{m}$ とする。$y$ は $x^2$ に比例し、速さが $40\mathrm{km/h}$ のときの制動距離は $8\mathrm{m}$ である。次の問いに答えなさい。}
\QQ{$y=ax^2$ と表すとき、比例定数 $a$ を求めなさい。}{}
\QQ{$y$ を $x$ の式で表しなさい。}{}
\QQ{速さが $60\mathrm{km/h}$ のときの制動距離を求めなさい。}{}
\QQ{速さを $40\mathrm{km/h}$ から $80\mathrm{km/h}$ にすると、制動距離は何倍になるか求めなさい。}{}
\TextTall{運転者が危険を感じてからブレーキがきき始めるまでに進む距離（空走距離）を $z\mathrm{m}$ とする。$z$ は $x$ に比例し、速さが $40\mathrm{km/h}$ のときの空走距離は $10\mathrm{m}$ である。また、空走距離と制動距離の和を停止距離という。}
\QQ{$z$ を $x$ の式で表しなさい。}{}
\QQ{停止距離を $w\mathrm{m}$ とするとき、$w$ を $x$ の式で表しなさい。}{}
\QQ{停止距離が $52\mathrm{m}$ であるとき、自動車の速さを求めなさい。}{}

% 振り子
\TypeQuestion{}{everyday_pendulum}
\TextTall{ある振り子が1往復するのにかかる時間を $x$ 秒、振り子の長さを $y\mathrm{cm}$ とすると、$y$ は $x^2$ に比例する。この振り子が1往復するのに $2$ 秒かかるとき、振り子の長さは $100\mathrm{cm}$ であった。次の問いに答えなさい。}
\QQ{$y$ を $x$ の式で表しなさい。}{}
\QQ{1往復するのに $1.6$ 秒かかる振り子の長さを求めなさい。}{}
\QQ{長さが $225\mathrm{cm}$ の振り子が1往復するのにかかる時間を求めなさい。}{}
\QQ{1往復するのにかかる時間を2倍にすると、振り子の長さは何倍になるか求めなさい。}{}
\QQ{振り子Aと振り子Bが1往復するのにかかる時間がそれぞれ $1.2$ 秒、$1.8$ 秒である。振り子Aと振り子Bの長さの比を、最も簡単な整数の比で表しなさい。}{}
\QQ{$0\leqq x\leqq 4$ の範囲で、$x$ と $y$ の関係をグラフに表しなさい。}{}

\LessonPageBreak

% タクシー料金：階段型
\TypeQuestion{}{everyday_taxi_fare}
\TextTall{あるタクシーの料金は、走行距離が $1.2\mathrm{km}$ までは $600$ 円であり、$1.2\mathrm{km}$ をこえると、その後は $300\mathrm{m}$ までごとに $100$ 円ずつ加算される。たとえば、走行距離が $1.2\mathrm{km}$ をこえて $1.5\mathrm{km}$ までの料金は $700$ 円である。次の問いに答えなさい。}
\QQ{走行距離が $1.2\mathrm{km}$ のときと $1.21\mathrm{km}$ のときの料金を、それぞれ求めなさい。}{}
\QQ{走行距離が $1.5\mathrm{km}$ のときと $1.51\mathrm{km}$ のときの料金を、それぞれ求めなさい。}{}
\QQ{走行距離が $2.05\mathrm{km}$ のときの料金を求めなさい。}{}
\QQ{料金が $800$ 円となる走行距離 $x\mathrm{km}$ の範囲を、不等号を用いて表しなさい。}{}
\QQ{$900$ 円で走ることのできる最長の距離を求めなさい。}{}
\QQ{走行距離を $x\mathrm{km}$、料金を $y$ 円として、$0<x\leqq 3$ の範囲で $x$ と $y$ の関係をグラフに表しなさい。}{}

% 電話料金：階段型
\TypeQuestion{}{everyday_phone_charge}
\TextTall{ある電話サービスの1か月の料金は、通話時間が $20$ 分までは $1200$ 円であり、$20$ 分をこえると、その後は $5$ 分までごとに $80$ 円ずつ加算される。たとえば、通話時間が $20$ 分をこえて $25$ 分までの料金は $1280$ 円である。次の問いに答えなさい。}
\QQ{通話時間が $20$ 分のときと $21$ 分のときの料金を、それぞれ求めなさい。}{}
\QQ{通話時間が $25$ 分のときと $26$ 分のときの料金を、それぞれ求めなさい。}{}
\QQ{通話時間が $47$ 分のときの料金を求めなさい。}{}
\QQ{料金が $1520$ 円となる通話時間 $x$ 分の範囲を、不等号を用いて表しなさい。}{}
\QQ{料金が $1600$ 円以下となる通話時間は、最長で何分か求めなさい。}{}
\QQ{通話時間を $x$ 分、料金を $y$ 円として、$0<x\leqq 60$ の範囲で $x$ と $y$ の関係をグラフに表しなさい。}{}
